\documentclass[header={Lecture 10 Exam Study Notes}]{study-note}

\begin{document}

\StudyTitleBlock{10}{Lecture 10: Exam Study Notes}{Spectrum, resolvent, Neumann series, and spectral radius}

\vspace{0.8em}

\begin{focusbox}
This lecture introduces spectral theory in Banach spaces. The exam usually expects the main definitions, the Neumann-series criterion, the openness of the resolvent set, and the core structural properties of the spectrum.
\end{focusbox}

\tableofcontents

\section{Definitions}

\begin{defbox}{Invertible operator}
For $T\in\mathcal L(X)$, the operator $T$ is \textbf{invertible} if there exists $T^{-1}\in\mathcal L(X)$ such that
\[
TT^{-1}=T^{-1}T=I.
\]
\end{defbox}

\begin{defbox}{Resolvent set and spectrum}
For a bounded operator $T\in\mathcal L(X)$ over a complex Banach space,
\[
\rho(T):=\{\lambda\in\C : \lambda I-T \text{ is invertible}\},
\qquad
\sigma(T):=\C\setminus \rho(T).
\]
The \textbf{point spectrum} is the set of eigenvalues.
\end{defbox}

\begin{warningbox}{Common trap}
For compact operators, only nonzero spectral values must be eigenvalues. The point $0$ is exceptional.
\end{warningbox}

\begin{warningbox}{Common trap}
The standard nonemptiness theorem for the spectrum requires complex scalars. Over real Banach spaces, one usually complexifies.
\end{warningbox}

\section{Neumann Series}

\begin{thmbox}{Theorem 14.1: Neumann series}
\thmNeumannSeries
\end{thmbox}

\begin{prooflevelbox}{FULL}
Know the geometric-series argument: $(I-T)^{-1}=\sum_{n=0}^\infty T^n$ when $\norm{T}<1$. Verify convergence, check it is a two-sided inverse, and state the norm estimate.
\end{prooflevelbox}

\paragraph{Why it matters.}
This is the basic perturbative invertibility criterion and appears directly in an oral-exam question.

\section{Resolvent And Spectrum}

\begin{thmbox}{Theorem 14.2: The resolvent set is open}
\thmResolventAnalytic
\end{thmbox}

\begin{prooflevelbox}{SKETCH}
Write $\lambda I-T=(\lambda_0 I-T)-(\lambda_0-\lambda)I$ and use the Neumann series to invert when $\abs{\lambda-\lambda_0}$ is small.
\end{prooflevelbox}

\begin{thmbox}{Resolvent identity}
If $\lambda,\mu\in\rho(T)$, then
\[
R(\lambda,T)-R(\mu,T)=(\mu-\lambda)R(\lambda,T)R(\mu,T).
\]
This is one of the standard algebraic identities used throughout spectral theory.
\end{thmbox}

\begin{thmbox}{Theorem 14.3: Basic spectral properties}
\thmSpectralRadius
\end{thmbox}

\begin{prooflevelbox}{SKETCH}
Know the three claims: $\sigma(T)$ is nonempty (Liouville theorem argument), compact, and $r(T)=\lim\norm{T^n}^{1/n}$. The spectral radius formula is the deepest part.
\end{prooflevelbox}

\section{Types Of Spectrum}

\begin{defbox}{Three pieces of the spectrum}
For $\lambda\in\sigma(T)$:
\begin{itemize}
\item \textbf{point spectrum}: $\lambda I-T$ is not injective;
\item \textbf{continuous spectrum}: $\lambda I-T$ is injective and has dense range, but is not surjective;
\item \textbf{residual spectrum}: $\lambda I-T$ is injective, but its range is not dense.
\end{itemize}
\end{defbox}

\begin{exbox}{What to keep in mind}
In a general Banach space, not every spectral value is an eigenvalue. This is one of the main differences from finite-dimensional linear algebra.
\end{exbox}

\section{Connection With Compact Operators}

\begin{focusbox}
Lecture 9 and Lecture 10 fit together as follows: if $K$ is compact and $\lambda\neq 0$, then studying $\lambda I-K$ is the same as studying
\[
I-\frac{1}{\lambda}K.
\]
Fredholm theory then shows that nonzero spectral values of compact operators are much more rigid.
\end{focusbox}

\section{What To Memorize Precisely}

\begin{focusbox}
Know precisely:
\begin{enumerate}
\item the definitions of invertibility, resolvent set, spectrum, and the three parts of the spectrum;
\item Theorems \textbf{14.1}, \textbf{14.2}, \textbf{14.3};
\item why Neumann series gives invertibility of $I-T$ when $\norm{T}<1$;
\item the resolvent identity;
\item that the spectrum of a bounded operator on a nontrivial complex Banach space is nonempty, compact, and contained in the closed disk of radius $\norm{T}$.
\end{enumerate}
\end{focusbox}

\end{document}
